% =================================================
% Решения домашней работы. Урок 3
% =================================================

\clearpage
\homeworktitle{Решения и комментарии}{Версия для учителя}

\begin{teacheranswer}[Задание 1]
  Сначала фиксируем цифру десятков:
  \[
    2:\ 22,26;\qquad 6:\ 62,66.
  \]
  Ответ:
  \[
    \boxed{22,\ 26,\ 62,\ 66}.
  \]
\end{teacheranswer}

\begin{criterionbox}
  Наличие чисел \(22\) и \(66\) показывает, что ученик правильно понял
  условие «цифры могут повторяться». Пропуск одного из них обычно означает,
  что ученик самовольно добавил запрет повторений.
\end{criterionbox}

\begin{teacheranswer}[Задание 2]
  Ноль не может стоять в разряде сотен. Поэтому начинаем с \(4\) или \(9\):
  \[
    4:\ 409,490;\qquad 9:\ 904,940.
  \]
  Ответ:
  \[
    \boxed{409,\ 490,\ 904,\ 940}.
  \]
\end{teacheranswer}

\begin{criterionbox}
  Записи \(049\) и \(094\) не являются трёхзначными числами: это числа
  \(49\) и \(94\).
\end{criterionbox}

\begin{teacheranswer}[Задание 3]
  Для наибольшего числа располагаем цифры по убыванию:
  \[
    \boxed{8530}.
  \]
  Для наименьшего сначала ставим наименьшую ненулевую цифру \(3\), затем
  ноль и оставшиеся цифры по возрастанию:
  \[
    \boxed{3058}.
  \]
\end{teacheranswer}

\begin{teacheranswer}[Задание 4]
  Если цифра единиц равна \(0,1,2,3,4,5,6,7\), то цифра десятков равна
  соответственно \(2,3,4,5,6,7,8,9\).
  \[
    \boxed{20,\ 31,\ 42,\ 53,\ 64,\ 75,\ 86,\ 97}.
  \]
\end{teacheranswer}

\begin{teacheranswer}[Задание 5]
  Последовательно меняем цифру десятков от \(1\) до \(9\):
  \[
    \boxed{18,\ 27,\ 36,\ 45,\ 54,\ 63,\ 72,\ 81,\ 90}.
  \]
  Всего \(9\) чисел.
\end{teacheranswer}

\begin{teacheranswer}[Задание 6]
  Группируем по цифре сотен:
  \[
    1:\ 147,174;
  \]
  \[
    4:\ 417,471;
  \]
  \[
    7:\ 714,741.
  \]
  Ответ: \(6\) чисел.
\end{teacheranswer}

\begin{criterionbox}
  Упорядоченный список позволяет увидеть структуру \(3\cdot2=6\): три
  варианта цифры сотен и два порядка оставшихся цифр.
\end{criterionbox}

\begin{teacheranswer}[Задание 7*]
  В разряд сотен можно поставить \(2\), \(5\) или \(7\): всего \(3\)
  варианта. После выбора сотен остаются \(3\) возможные цифры для десятков,
  а затем \(2\) цифры для единиц.
  \[
    3\cdot3\cdot2=\boxed{18}.
  \]
\end{teacheranswer}

\begin{criterionbox}
  Важно, чтобы ученик объяснил, почему для первой позиции не четыре, а три
  варианта: ноль не может быть цифрой сотен трёхзначного числа.
\end{criterionbox}

\begin{teacheranswer}[Задание 8]
  \[
    6\text{ десятков тысяч}+128\text{ сотен}
    =6\cdot10\,000+128\cdot100
  \]
  \[
    =60\,000+12\,800=\boxed{72\,800}.
  \]
\end{teacheranswer}
